\section*{अध्याय \ref{ch:g8:reproductive-systems} --- जनन-तंत्र}

\begin{solution}{exo:g8:reproductive-systems:1}
\omterm{def:g8:reproductive-systems:male}{वृषण}, वृषणकोश में: \omterm{def:g8:puberty:puberty}{यौवनारंभ} से लगातार शुक्राणुओं का बनना। नलियाँ और
ग्रंथियाँ: जमा करना, ढोना, और वे तरल जो वीर्य बनाते हैं। \omterm{def:g8:reproductive-systems:male}{शिश्न}: वीर्य का
निकास --- और अपनी अलग, \omterm{def:g7:heart-and-circulation:rooms}{कपाटों} से सँभली ड्यूटी पर \omterm{def:g7:removing-wastes:kidneys}{मूत्र} का भी।
\end{solution}

\begin{solution}{exo:g8:reproductive-systems:2}
\omterm{def:g8:reproductive-systems:female}{अंडाशय}: जन्म से \omterm{def:g4:animal-reproduction:sexual}{अंडाणुओं} का भंडार, और \omterm{def:g8:puberty:puberty}{यौवनारंभ} से हर चक्र में एक छोड़ता
हुआ। \omterm{def:g8:reproductive-systems:female}{अंडवाहिनियाँ}: छोड़े गए \omterm{def:g2:animals-grow-up:born}{अंडे} को लपकती हुईं --- यानी \omterm{def:g5:flower-to-fruit:steps}{निषेचन} की आम
जगह। \omterm{prop:g5:human-reproduction:plan}{गर्भाशय}: पेशीदार, और नई होती परत वाला ठिकाना। योनि: वीर्य लेने की
जगह, और जन्म की नली।
\end{solution}

\begin{solution}{exo:g8:reproductive-systems:3}
\omterm{def:g8:reproductive-systems:male}{शुक्राणु}: शरीर की सबसे छोटी \omterm{def:g5:first-look-at-cells:cell}{कोशिका} --- \omterm{def:g6:cell-unit-of-life:parts}{केंद्रक} से भरा एक \omterm{def:g1:my-body:parts}{सिर} और तैरने
वाली एक पूँछ, हलका-सा साज़-सामान लिए, और लगातार करोड़ों की तादाद में
बना हुआ। \omterm{def:g4:animal-reproduction:sexual}{अंडाणु}: शरीर की सबसे बड़ी, गोल और बिना हिली-डुली, भरपूर
साज़-सामान वाली, और महीने में एक छोड़ी हुई।
\end{solution}

\begin{solution}{exo:g8:reproductive-systems:4}
कोई \omterm{def:g8:reproductive-systems:female}{अंडाशय} एक \omterm{def:g4:animal-reproduction:sexual}{अंडाणु} छोड़ता है, जिसे उसकी \omterm{def:g8:reproductive-systems:female}{अंडवाहिनी} लपक लेती है ---
और यह चौदहवें दिन के आस-पास होता है। वह लगभग एक दिन तक निषेचित हो सकता
है।
\end{solution}

\begin{solution}{exo:g8:reproductive-systems:5}
\omterm{prop:g5:human-reproduction:plan}{गर्भाशय} की तैयार परत का टूटना और गिरना, जो कुछ दिनों में बाहर चली जाती
है। उसका आना कहता है कि वह चक्र बिना \omterm{def:g5:flower-to-fruit:steps}{निषेचन} के ख़त्म हुआ।
\end{solution}

\begin{solution}{exo:g8:reproductive-systems:6}
\omterm{def:g5:flower-to-fruit:steps}{निषेचन}। \omterm{prop:g8:sexual-reproduction:core}{युग्मनज} तैयार परत में जम जाता है, चक्र ठहर जाता है, और
\omterm{def:g5:human-reproduction:pregnancy}{गर्भावस्था} शुरू हो जाती है।
\end{solution}

\begin{solution}{exo:g8:reproductive-systems:7}
पहला क़दम, \omterm{def:g8:sexual-reproduction:gametes}{युग्मक}: \omterm{def:g8:reproductive-systems:male}{वृषण} और \omterm{def:g8:reproductive-systems:female}{अंडाशय}। दूसरा क़दम, \omterm{def:g5:flower-to-fruit:steps}{निषेचन}: \omterm{def:g8:reproductive-systems:female}{अंडवाहिनी} में,
जहाँ \omterm{def:g8:reproductive-systems:male}{शुक्राणु} योनि से \omterm{prop:g5:human-reproduction:plan}{गर्भाशय} होते हुए तैरकर पहुँचते हैं। तीसरा क़दम,
बँटते हुए \omterm{prop:g8:sexual-reproduction:core}{युग्मनज} का विकास: \omterm{prop:g5:human-reproduction:plan}{गर्भाशय} की तैयार परत में ठहरा हुआ।
\end{solution}

\begin{solution}{exo:g8:reproductive-systems:8}
\omterm{def:g4:animal-reproduction:sexual}{नर} पक्ष ज़्यादा-और-सस्ता वाला तर्क चलाता है: लगातार तैरते लाखों, और हर
एक ख़र्च करने लायक़। \omterm{def:g4:animal-reproduction:sexual}{मादा} पक्ष कम-और-साज़-सामान वाला चलाता है: महीने में
एक भरा-पूरा \omterm{def:g4:animal-reproduction:sexual}{अंडाणु}, और उसके पीछे \omterm{prop:g5:human-reproduction:plan}{गर्भाशय} के ठिकाने का पूरा निवेश। यानी
रणनीति की वही लकीर, \omterm{def:g5:first-look-at-cells:cell}{कोशिकाओं} में लिखी हुई।
\end{solution}

\begin{solution}{exo:g8:reproductive-systems:9}
पहले कुछ सालों में मशीनरी सध रही होती है --- \omterm{def:g8:puberty:hormones}{हॉर्मोन} की कड़ी अपनी लय में
बैठ रही होती है --- इसलिए अनियमितता अपेक्षित है; और चक्र की लंबाई सचमुच
बदलती है, व्यक्ति-दर-व्यक्ति और महीने-दर-महीने। दोनों ही तंत्र के
सामान्य ढंग से चलने की निशानी हैं।
\end{solution}

\begin{solution}{exo:g8:reproductive-systems:10}
\omterm{prop:g5:human-reproduction:plan}{गर्भाशय} एक \omterm{def:g3:bones-and-muscles:muscle}{पेशी} है, और परत का गिरना \omterm{def:g3:bones-and-muscles:muscle}{पेशी} का काम है --- उसके \omterm{def:g3:bones-and-muscles:muscle}{संकुचन} ही
ऐंठन के रूप में महसूस होते हैं। गरमाहट और हरकत मदद करती हैं (और जब वे
काफ़ी न हों तो स्कूल की नर्स के पास सलाह है)।
\end{solution}

\begin{solution}{exo:g8:reproductive-systems:11}
\omterm{def:g8:reproductive-systems:male}{शुक्राणु} का काम सफ़र है: सस्ता, तैरने वाला, और लाखों में चाहिए क्योंकि
पहुँचता क़रीब-क़रीब कोई नहीं --- इसलिए कारख़ाना लगातार चलता है। और \omterm{def:g2:animals-grow-up:born}{अंडे}
का काम \omterm{prop:g8:sexual-reproduction:core}{युग्मनज} का पहला साज़-सामान बनना है: महँगा, जन्म से भंडार में
रखा, और अपने पीछे एक महीने की तैयारी लेकर छोड़ा हुआ --- इसलिए वह
कारख़ाना एक ख़ज़ाना है, जो एक बार में एक ही ख़र्च करता है।
\end{solution}

\begin{solution}{exo:g8:reproductive-systems:12}
पहली घटना वह सवाल तैयार कर रही है: परत दोबारा बनी, ठिकाना तैयार। दूसरी
घटना सवाल पूछती है: एक \omterm{def:g4:animal-reproduction:sexual}{अंडाणु} छोड़ा गया --- क्या \omterm{def:g5:flower-to-fruit:steps}{निषेचन} आएगा? तीसरी घटना
जवाब है ``नहीं'': बिना ज़रूरत की परत गिर जाती है, और पहिया अगले महीने
फिर पूछता है। और चौथी घटना जवाब है ``हाँ'': \omterm{prop:g8:sexual-reproduction:core}{युग्मनज} जम जाता है, पूछना
रुक जाता है, और वह ठिकाना नौ महीने अपना मक़सद पूरा करता है।
\end{solution}

\begin{solution}{pb:g8:reproductive-systems:1}
\textbf{1.} \omterm{def:g8:reproductive-systems:male}{वृषणों} से --- यानी उन कारख़ानों से, जो वृषणकोश में रखे हैं
--- \omterm{def:g8:reproductive-systems:male}{शुक्राणु} जमा होने तथा ढोए जाने के लिए नलियों में जाते हैं, रास्ते
में ग्रंथियों के तरल उनसे जुड़कर वीर्य बनाते हैं, और वे \omterm{def:g8:reproductive-systems:male}{शिश्न} से बाहर
जाते हैं; वही अंग बाक़ी वक़्त \omterm{def:g7:removing-wastes:kidneys}{मूत्र} ले जाता है, और \omterm{def:g7:heart-and-circulation:rooms}{कपाटों} का एक इंतज़ाम
दोनों ड्यूटियों को सख़्ती से अलग रखता है।

\textbf{2.} भंडार \omterm{def:g8:reproductive-systems:female}{अंडाशयों} में इंतज़ार करता है, जन्म से ही; \omterm{def:g5:flower-to-fruit:steps}{निषेचन} आम
तौर पर \omterm{def:g8:reproductive-systems:female}{अंडवाहिनियों} में होता है; और \omterm{def:g5:human-reproduction:pregnancy}{गर्भावस्था} \omterm{prop:g5:human-reproduction:plan}{गर्भाशय} की तैयार परत में
ठहरती है।

\textbf{3.} दोनों \omterm{def:g8:sexual-reproduction:gametes}{युग्मक} उलटी रणनीतियों पर बने हैं: लाखों सस्ते तैराक,
बनाम महीने में एक साज़-सामान वाला \omterm{def:g2:animals-grow-up:born}{अंडा} --- यानी संख्या-बनाम-देखभाल,
ख़ुद \omterm{def:g5:first-look-at-cells:cell}{कोशिकाओं} में लिखा हुआ।

\textbf{4.} अपने हुक्म का इंतज़ार कर रहे थे: तंत्र बने-बनाए तो थे, पर
बेकार पड़े थे, जब तक \omterm{def:g5:brain-in-command:system}{मस्तिष्क} के तल वाले कमान-केंद्र के हॉर्मोनों ने
उन्हें \omterm{def:g8:puberty:puberty}{यौवनारंभ} पर चालू नहीं कर दिया।

\textbf{5.} पहले से लेकर पाँचवें दिन के आस-पास: \omterm{prop:g8:reproductive-systems:cycle}{मासिक धर्म}। चक्र के बीच
तक: परत दोबारा बनती है। चौदहवें दिन के आस-पास: \omterm{prop:g8:reproductive-systems:cycle}{अंडोत्सर्ग}, और उसके
इर्द-गिर्द उपजाऊ दिन। फिर तैयार रहकर इंतज़ार --- अट्ठाईसवें दिन और
पहिए के अगले चक्कर तक।

\textbf{6.} चक्र की लंबाई व्यक्ति-दर-व्यक्ति बदलती है, और एक ही व्यक्ति
में महीने-दर-महीने भी --- और पहले कुछ साल तो मशीनरी सधने तक अनियमित ही
चलते हैं। कैलेंडर ढर्रा सिखाता है, किसी की ठीक-ठीक तारीख़ें नहीं।

\textbf{7.} रद्द: परत का टूटना और अगला \omterm{prop:g8:reproductive-systems:cycle}{मासिक धर्म} --- पहिया रुक जाता
है। और शुरू: \omterm{prop:g8:sexual-reproduction:core}{युग्मनज} का परत में जमना, तथा \omterm{def:g5:human-reproduction:pregnancy}{गर्भावस्था} के नौ महीने।

\textbf{8.} ऐंठन पेशीदार \omterm{prop:g5:human-reproduction:plan}{गर्भाशय} का परत गिराते हुए काम करना है ---
सामान्य, कोई नुक़सान नहीं; गरमाहट और हरकत मदद करती हैं, और उससे कड़ी हर
बात का पता नर्स या डॉक्टर का दरवाज़ा है।

\textbf{9.} मिसाल के लिए: ``ये अंग भी \omterm{def:g4:heart-and-blood:heart}{दिल} और फेफड़ों जैसे ही अंग हैं ---
इस कमरे में मौजूद हर किसी के पास इनमें से एक या दूसरा जोड़ा है, और यह
बस उनका अध्याय है।''

\textbf{10.} कारख़ाने: \omterm{def:g8:reproductive-systems:male}{वृषण} --- \omterm{def:g8:reproductive-systems:female}{अंडाशय}। \omterm{def:g8:sexual-reproduction:gametes}{युग्मक}: लगातार तैरते लाखों ---
या महीने में एक साज़-सामान वाला \omterm{def:g2:animals-grow-up:born}{अंडा}। मिलन: \omterm{def:g8:reproductive-systems:female}{अंडवाहिनी}। ठिकाना: कोई नहीं
--- या हर महीने तैयार होता \omterm{prop:g5:human-reproduction:plan}{गर्भाशय}। समय-सारणी: \omterm{def:g8:puberty:puberty}{यौवनारंभ} से लगातार ---
या चक्रीय, \omterm{def:g8:puberty:puberty}{यौवनारंभ} से लगभग पचास तक।

\textbf{11.} दोनों नमूनों पर पहला क़दम: कारख़ाने \omterm{def:g8:sexual-reproduction:gametes}{युग्मक} बनाते हैं। \omterm{def:g4:animal-reproduction:sexual}{मादा}
नमूने पर दूसरा क़दम: उनका \omterm{def:g8:reproductive-systems:female}{अंडवाहिनी} में जुड़ना। और नमूने तथा कैलेंडर,
दोनों पर तीसरा क़दम: \omterm{prop:g8:sexual-reproduction:core}{युग्मनज} का कैलेंडर वाली तैयार परत में ठहरना, पहिए
का रुकना, और विकास का शुरू होना।

\textbf{12.} ``क्योंकि शरीर गप के मुक़ाबले जानकारी पर बेहतर चलते हैं ---
और आज सिखाई हर बात के आगे के सवालों का एक प्रशिक्षित, गोपनीय पता भी
है।''
\end{solution}
